\documentclass[11pt,reqno]{amsart}

\usepackage{lmodern}
\usepackage[T1]{fontenc}
\usepackage[utf8]{inputenc}
\usepackage[english]{babel}
\usepackage{microtype}
\usepackage{amsmath,amssymb,amsthm,mathtools}
\usepackage{booktabs}
\usepackage{enumitem}
\usepackage{xcolor}
\usepackage{hyperref}
\usepackage{ytableau}
\ytableausetup{centertableaux,boxsize=1.4em}

% ---------- Theorems ----------
\newtheorem{theorem}{Theorem}[section]
\newtheorem{proposition}[theorem]{Proposition}
\newtheorem{lemma}[theorem]{Lemma}
\newtheorem{corollary}[theorem]{Corollary}
\newtheorem{conjecture}[theorem]{Conjecture}
\theoremstyle{definition}
\newtheorem{definition}[theorem]{Definition}
\newtheorem{example}[theorem]{Example}
\newtheorem{remark}[theorem]{Remark}
\newtheorem{observation}[theorem]{Observation}
\newtheorem{openproblem}[theorem]{Open problem}

% ---------- Macros ----------
\newcommand{\setN}{\mathbb{N}}
\newcommand{\setZ}{\mathbb{Z}}
\newcommand{\setR}{\mathbb{R}}

\newcommand{\defin}[1]{%
\relax\ifmmode%
\textcolor{blue}{#1}%
\else \textcolor{blue}{\emph{#1}}%
\fi%
}

\DeclareMathOperator{\rk}{rk}

\newcommand{\interl}{\ll}

\title{Towards real-rootedness of non-nesting rook polynomials}
\author{Per Alexandersson}
\address{Department of Mathematics, Stockholm University}
\email{per.alexandersson@math.su.se}
\date{\today}

\begin{document}

\begin{abstract}
We study the generating polynomial for non-nesting rook placements
on Ferrers boards, weighted by the number of rooks.
We conjecture that these polynomials are always real-rooted.
We record a recursive interlacing framework, prove the conjecture in
the staircase and rectangular cases, and isolate a difference-quotient
recurrence package which would imply the full result.
Computations for all partitions with $|\mu| \leq 20$
support the missing interlacing statement.
Special cases recover the Narayana polynomials (staircase shapes)
and squared binomial coefficients (rectangular shapes).
\end{abstract}

\maketitle

\section{Introduction}

Let $\mu = (\mu_1 \geq \mu_2 \geq \dotsb \geq \mu_\ell > 0)$
be a partition, defining a Ferrers board in English convention.
A \defin{rook placement} on $\mu$ is a set of cells, at most one per
row and at most one per column.
A rook placement is \defin{non-nesting} if no rook is strictly
north-west of another; equivalently, reading rooks left to right
by column, the row indices are strictly decreasing.

\begin{definition}
The \defin{non-nesting rook polynomial} of $\mu$ is
\[
   \defin{R_\mu(t)} \coloneqq \sum_{k \geq 0}
   r_k(\mu) \, t^k,
\]
where $r_k(\mu)$ is the number of non-nesting rook placements
on $\mu$ with exactly $k$ rooks.
\end{definition}

\begin{example}\label{ex:small}
For the staircase $\delta_3 = (3,2,1)$, we have $R_{\delta_3}(t) = 1 + 6t + 6t^2 + t^3$.
For the $3 \times 3$ square, $R_{(3,3,3)}(t) = 1 + 9t + 9t^2 + t^3$.
\end{example}

Our main conjecture is:

\begin{conjecture}\label{conj:main}
For every partition $\mu$, the polynomial $R_\mu(t)$ is real-rooted.
\end{conjecture}

\section{Recursion and special cases}

The polynomial $R_\mu(t)$ satisfies the following recursion,
obtained by conditioning on whether a rook is placed in the
last (shortest) row, and if so, in which column.

\begin{figure}[ht]
\centering
\ytableausetup{boxsize=1.1em}
\[
R\Bigl(\;
\begin{ytableau}
{} & {} & {} \\
{} & {} \\
*(yellow!40)
\end{ytableau}
\;\Bigr)
\;=\;
R\Bigl(\;
\begin{ytableau}
{} & {} & {} \\
{} & {}
\end{ytableau}
\;\Bigr)
\;+\; t \cdot R\Bigl(\;
\begin{ytableau}
*(lightgray) & {} & {} \\
*(lightgray) & {}
\end{ytableau}
\;\Bigr)
\]
\ytableausetup{boxsize=1.4em}
\caption{The recursion for $\mu = (3,2,1)$.
The bottom row (yellow) is processed: either no rook is placed,
or a rook occupies column $c = 1$ of the last row
and columns $\leq c$ (gray) become forbidden for remaining rooks.}
\label{fig:recursion}
\end{figure}

\begin{lemma}\label{lem:recursion}
Let $\mu = (\mu_1, \dotsc, \mu_\ell)$ with $\ell \geq 2$,
and set $\mu' = (\mu_1, \dotsc, \mu_{\ell-1})$.
Then
\begin{equation}\label{eq:recursion}
   R_\mu(t)
   = R_{\mu'}(t)
   + t \sum_{c=1}^{\mu_\ell}
     R_{\mu'_c}(t),
\end{equation}
where $\mu'_c = (\max(0, \mu_1 - c), \dotsc, \max(0, \mu_{\ell-1} - c))$
with zero parts removed.
\end{lemma}
\begin{proof}
If no rook occupies the last row, the remaining placements
live on $\mu'$.
If a rook occupies column $c$ of the last row ($1 \leq c \leq \mu_\ell$),
then every remaining rook must use a column $> c$ (otherwise a rook
in a higher row at column $\leq c$ would be north-west of the bottom rook).
Relabeling columns $c+1, c+2, \dotsc$ as $1, 2, \dotsc$
gives a Ferrers board with parts $\max(0, \mu_i - c)$.
\end{proof}


\subsection{Staircase shapes}
For the staircase partition $\delta_n = (n, n{-}1, \dotsc, 1)$,
the polynomial $R_{\delta_n}(t)$ is the $n$th row of the
Narayana triangle (shifted by one), i.e.,
\[
   R_{\delta_n}(t) = \sum_{k=0}^{n} N(n{+}1, k{+}1) \, t^k
   = \sum_{k=0}^{n}
     \frac{1}{n{+}1}\binom{n{+}1}{k{+}1}\binom{n{+}1}{k} \, t^k.
\]
These are palindromic and real-rooted (known, e.g., from the theory of
non-crossing partitions).

\subsection{Rectangular shapes}
For the $n \times n$ square $\mu = (n^n)$,
\[
   R_{(n^n)}(t)
   = \sum_{k=0}^{n} \binom{n}{k}^2 t^k.
\]
In particular, $R_{(n^n)}(1) = \binom{2n}{n}$ and the polynomial
is palindromic. Real-rootedness follows from the fact that the
squared binomial coefficients $\binom{n}{k}^2$ are the coefficients of
$(1+t)^n \cdot (1+t)^n$ projected onto the diagonal, or equivalently
from Legendre polynomial theory.

More generally, for a rectangle $\mu = (m^n)$ with $m \geq n$,
\[
   R_{(m^n)}(t) = \sum_{k=0}^{n} \binom{m}{k}\binom{n}{k} \, t^k.
\]
In particular, $R_{(m^n)}(1) = \binom{m+n}{n}$ counts lattice paths, and
$R_{(n^n)}(1) = \binom{2n}{n}$.


\subsection{Table of small examples}

\begin{center}
\small
\begin{tabular}{lccc}
\toprule
$\mu$ & $R_\mu(t)$ & palindromic & gamma vector \\
\midrule
$(1)$       & $1+t$             & --- & --- \\
$(2)$       & $1+2t$            & --- & --- \\
$(1,1)$     & $1+2t$            &     & --- \\
$(2,1)$     & $1+3t+t^2$        & \checkmark & $(1,1)$ \\
$(2,2)$     & $1+4t+t^2$        & \checkmark & $(1,2)$ \\
$(3,1)$     & $1+4t+2t^2$       &     & --- \\
$(3,2,1)$   & $1+6t+6t^2+t^3$   & \checkmark & $(1,3)$ \\
$(3,3,1)$   & $1+7t+7t^2+t^3$   & \checkmark & $(1,4)$ \\
$(3,3,3)$   & $1+9t+9t^2+t^3$   & \checkmark & $(1,6)$ \\
$(4,3,2,1)$ & $1+10t+20t^2+10t^3+t^4$ & \checkmark & $(1,6,2)$ \\
\bottomrule
\end{tabular}
\end{center}


\section{A Br\"and\'en-style interlacing attempt}

A natural first approach to Conjecture~\ref{conj:main}
is to try to show that the row-by-row recursion preserves
an interlacing chain, via the framework of Br\"and\'en~\cite{Branden2015}.

\begin{definition}\label{def:interlacing}
Let $f$ and $g$ be polynomials with positive leading coefficients and
all real, non-positive roots.
We say $f$ \defin{(weakly) interlaces} $g$, written $f \interl g$,
if the roots alternate: $\dotsb \leq a_2 \leq b_2 \leq a_1 \leq b_1 \leq 0$.
A sequence $(f_1, \dotsc, f_n)$ of real-rooted polynomials is
\defin{interlacing} if $f_i \interl f_j$ for all $i < j$.
\end{definition}

We recall the key tool from Br\"and\'en's theory.

\begin{theorem}[{\cite[Cor.~8.7]{Branden2015}}]\label{thm:branden}
Let $T = (t_{ij})$ be an $m \times n$ matrix of linear forms
$t_{ij} = a_{ij} + b_{ij} x$ with $a_{ij}, b_{ij} \geq 0$.
Suppose that for each row~$i$, there exists an index $p_i$ such that
$b_{ij} > 0$ only for $j \leq p_i$ and
$a_{ij} > 0$ only for $j > p_i$,
and $p_1 \leq p_2 \leq \dotsb \leq p_m$.
If $(f_1, \dotsc, f_n)$ is an interlacing sequence,
then the nonzero entries of $T \cdot (f_1, \dotsc, f_n)^T$
form an interlacing sequence.
\end{theorem}

\begin{definition}
For a partition $\mu$ and integer $c \geq 0$, write
\[
  \mu_{\langle c \rangle} \coloneqq
  (\max(0, \mu_1 {-} c), \dotsc, \max(0, \mu_\ell {-} c))
  \quad\text{(zero parts removed)}
\]
for the board obtained by stripping the first $c$ columns from~$\mu$.
Note $\mu_{\langle 0 \rangle} = \mu$.
\end{definition}


\begin{figure}[ht]
\centering
\ytableausetup{boxsize=1.1em}
\begin{tabular}{c@{\quad}c@{\quad}c@{\quad}c@{\quad}c@{\quad}c@{\quad}c}
$\mu_{\langle 3 \rangle}$ && $\mu_{\langle 2 \rangle}$ &&
$\mu_{\langle 1 \rangle}$ && $\mu_{\langle 0 \rangle} = \mu$
\\[4pt]
$\varnothing$
& $\interl$ &
\begin{ytableau}
{}
\end{ytableau}
& $\interl$ &
\begin{ytableau}
{} & {} \\
{}
\end{ytableau}
& $\interl$ &
\begin{ytableau}
{} & {} & {} \\
{} & {} \\
{}
\end{ytableau}
\\[6pt]
$1$ && $1{+}t$ && $1{+}3t{+}t^2$ && $1{+}6t{+}6t^2{+}t^3$
\end{tabular}
\ytableausetup{boxsize=1.4em}

\medskip
\caption{Column-strip sequence for $\mu = (3,2,1)$.
Stripping $c$ columns from the left yields $\mu_{\langle c \rangle}$;
the non-nesting rook polynomials form an interlacing chain.}
\label{fig:column_strip}
\end{figure}

A natural strengthened conjecture is the following.

\begin{conjecture}\label{conj:main_strong}
For every partition $\mu$, the column-strip sequence
\[
  R_{\mu_{\langle M \rangle}} \interl
  R_{\mu_{\langle M-1 \rangle}} \interl
  \dotsb \interl
  R_{\mu_{\langle 1 \rangle}} \interl R_\mu
\]
is interlacing, where $M = \mu_1$.
In particular, $R_\mu$ is real-rooted.
\end{conjecture}

A first Br\"and\'en-style induction suggests why
Conjecture~\ref{conj:main_strong} might be true.
Let $m = \mu_\ell$, $\mu' = (\mu_1,\dotsc,\mu_{\ell-1})$, and
\begin{equation}\label{eq:V_chain}
  f_i \coloneqq R_{\mu'_{\langle m+1-i \rangle}}
  \qquad (1 \leq i \leq m+1).
\end{equation}
Then the recursion of Lemma~\ref{lem:recursion} gives
\[
  g_k
  \coloneqq R_{\mu_{\langle m+1-k \rangle}}
  = f_k + t(f_1+\dotsb+f_{k-1}),
  \qquad 1 \leq k \leq m+1.
\]
Equivalently, $(g_1,\dotsc,g_{m+1})^T = N (f_1,\dotsc,f_{m+1})^T$ where
\[
  N =
  \begin{pmatrix}
  1 & 0 & 0 & \cdots & 0 \\
  t & 1 & 0 & \cdots & 0 \\
  t & t & 1 & \cdots & 0 \\
  \vdots & & & \ddots & \vdots \\
  t & t & t & \cdots & 1
  \end{pmatrix}.
\]
If the sequence $(f_1,\dotsc,f_{m+1})$ were pairwise interlacing,
Br\"and\'en's Corollary~8.7 would imply that
$(g_1,\dotsc,g_{m+1})$ is interlacing as well.
The difficulty is that the outer induction supplies only
\emph{adjacent} interlacing for the strip sequence of~$\mu'$,
and that is not enough to invoke Corollary~8.7.
This is the gap in the naive proof.

\begin{remark}[Current Lean status]
The current Lean development verifies the matrix-side structure of the rook
update and proves termwise real-rootedness and nonnegativity for
\[
  g_k = f_k + t \cdot (f_1 + \dotsb + f_{k-1}).
\]
What has not yet been formalized there is the final pairwise interlacing-chain
step that the naive Br\"and\'en-style argument would need.
\end{remark}

\begin{remark}
The analogous result for the \emph{standard} rook polynomial
(nesting allowed) also holds: the polynomials are real-rooted
for all Ferrers boards.
For staircases, the standard rook polynomial gives
the reversed Stirling numbers of the second kind.
The standard recursion $R_\mu = R_{\mu'} + \mu_\ell \cdot t \cdot R_{\mu'-1}$
(where $\mu'-1$ subtracts~$1$ from each part) factors into
a product via the Goldman--Joichi--White theorem.
\end{remark}

\section{Touchard Polynomials and Big Blocks}\label{sec:touchard}

Let $\Pi_n$ denote the set of set partitions of $[n] = \{1,\dotsc,n\}$,
and let $\ell(\pi)$ be the number of blocks of~$\pi$.

\subsection{Touchard polynomials}

The \defin{Touchard polynomials}\footnote{\url{https://en.wikipedia.org/wiki/Touchard_polynomials}}
(also called Bell polynomials) are
\[
  B_n(x)
  \coloneqq
  \sum_{k=1}^{n} S(n,k)\,x^k
  =
  \sum_{\pi \in \Pi_n} x^{\ell(\pi)},
\]
where $S(n,k)$ is a Stirling number of the second kind.

The \defin{multivariate Bell polynomial} is
\[
  B_n(x_1,\dotsc,x_n)
  \coloneqq
  \sum_{\pi \in \Pi_n}
  \prod_{B \in \pi} x_{|B|}.
\]
For example,
\[
  B_3(x_1,x_2,x_3) = x_1^3 + 3x_1x_2 + x_3,
\]
corresponding to the set partitions
$1|2|3$, $1|23$, $12|3$, $13|2$, and $123$.

This multivariate refinement is not same-phase stable in general.
Indeed,
\[
  B_3(3t,t,t) = t + 9t^2 + 27t^3,
\]
which is not real-rooted.

The reason this family is relevant here is that the staircase
\emph{standard} rook polynomial is the reversal of a Touchard polynomial.
If $R_\mu^{\mathrm{std}}(t)$ denotes the standard rook polynomial
(nesting allowed), then for the staircase $\delta_{n-1} = (n-1,n-2,\dotsc,1)$,
\[
  R_{\delta_{n-1}}^{\mathrm{std}}(t)
  =
  \sum_{k=0}^{n-1} S(n,n-k)\,t^k
  =
  t^n B_n(t^{-1}).
\]
Thus the standard staircase-rook family is a rook model for the
Touchard polynomials, up to coefficient reversal.

\begin{openproblem}\label{open:touchard_multivar}
Find a multivariate analog of $B_n(x)$ that is stable, or at least
same-phase stable, and ideally is multi-affine and admits a natural
recursive description compatible with the staircase-rook model.
\end{openproblem}

\subsection{Counting big blocks}

For $j \geq 2$, define
\[
  \operatorname{bb}_j(\pi)
  \coloneqq
  \#\{B \in \pi : |B| \geq j\},
\]
the number of blocks of size at least~$j$, and set
\[
  P_{n,j}(t)
  \coloneqq
  \sum_{\pi \in \Pi_n} t^{\operatorname{bb}_j(\pi)}.
\]

\medskip

For the staircase board it is convenient to reverse the columns and work on the
triangular board
\[
  T_n
  \coloneqq
  \{(i,j) : 1 \leq i < j \leq n\}.
\]
Since standard rook numbers are unchanged by permuting columns, this is
equivalent to working on~$\delta_{n-1}$ itself. If $\rho$ is a standard rook
placement on~$T_n$, let $D(\rho)$ be the directed graph on~$[n]$ with an edge
$i \to j$ for each rook~$(i,j)$, and define
\[
  \operatorname{rch}_j(\rho)
  \coloneqq
  \#\{\text{path components of $D(\rho)$ having at least $j$ vertices}\}.
\]
Equivalently, $\operatorname{rch}_j(\rho)$ counts maximal rook chains
\[
  (a_1,a_2),\ (a_2,a_3),\ \dotsc,\ (a_{m-1},a_m)
\]
with $m \geq j$.

\begin{proposition}\label{prop:touchard_rook_bijection}
The map sending a standard rook placement $\rho$ on~$T_n$ to the collection of
vertex sets of the path components of~$D(\rho)$ is a bijection from staircase
rook placements to set partitions of~$[n]$. Under this bijection, a block of
size~$m$ corresponds to a path component with $m$ vertices, equivalently to a
maximal rook chain with $m-1$ rooks. Consequently,
\[
  P_{n,j}(t)
  =
  \sum_{\rho} t^{\operatorname{rch}_j(\rho)},
\]
where the sum ranges over all standard rook placements $\rho$ on~$T_n$.
In particular,
\[
  R_{\delta_{n-1}}^{\mathrm{std}}(t)
  =
  \sum_{\rho} t^{|\rho|}
  =
  \sum_{\pi \in \Pi_n} t^{n-\ell(\pi)}
  =
  t^n B_n(t^{-1}).
\]
\end{proposition}
\begin{proof}
Because $\rho$ has at most one rook in each row and column, every vertex of
$D(\rho)$ has outdegree at most~$1$ and indegree at most~$1$. Since each edge
has the form $i \to j$ with $i<j$, the graph is acyclic, hence every connected
component is a directed path. The vertex sets of these path components
therefore form a set partition of~$[n]$.

Conversely, if $B = \{b_1 < \dotsb < b_m\}$ is a block of a set partition,
place rooks at
\[
  (b_1,b_2),\ (b_2,b_3),\ \dotsc,\ (b_{m-1},b_m).
\]
Doing this for each block produces a standard rook placement on~$T_n$, since
different blocks give disjoint sets of left endpoints and right endpoints.
These two constructions are inverse to each other.

Finally, a block of size at least~$j$ corresponds exactly to a path component
with at least~$j$ vertices, or equivalently to a maximal rook chain of length
at least~$j-1$. This proves the formula for~$P_{n,j}(t)$. The displayed
formula for $R_{\delta_{n-1}}^{\mathrm{std}}(t)$ is the special case obtained
by weighting a placement by its number of rooks.
\end{proof}

\begin{proposition}\label{prop:big_block_recur}
For $n \geq 0$ and $j \geq 2$,
\begin{equation}\label{eq:big_block_recur}
  P_{n+1,j}(t)
  =
  \sum_{k=0}^{j-2} \binom{n}{k} P_{n-k,j}(t)
  + t \sum_{k=j-1}^{n} \binom{n}{k} P_{n-k,j}(t).
\end{equation}
\end{proposition}
\begin{proof}
Consider the block containing $n+1$.
If it has size $k+1$, then one chooses its other $k$ elements in
$\binom{n}{k}$ ways and partitions the remaining $n-k$ elements
arbitrarily.
If $k \leq j-2$, then the new block is not counted by
$\operatorname{bb}_j$; if $k \geq j-1$, then it contributes one extra
big block, hence a factor of~$t$.
Summing over all $k$ gives~\eqref{eq:big_block_recur}.
\end{proof}

Equivalently, the exponential generating function is
\begin{equation}\label{eq:big_block_egf}
  \sum_{n \geq 0} P_{n,j}(t)\,\frac{z^n}{n!}
  =
  \exp\!\left(
    \sum_{r=1}^{j-1} \frac{z^r}{r!}
    +
    t \sum_{r=j}^{\infty} \frac{z^r}{r!}
  \right).
\end{equation}

Now define
\[
  Q_{n,j}(t) \coloneqq t\,P_{n,j}(t).
\]

\begin{proposition}\label{prop:Q_recur}
For $n \geq 0$ and $j \geq 2$,
\begin{equation}\label{eq:Q_recur}
  Q_{n+1,j}(t)
  =
  t\,Q'_{n,j}(t)
  +
  \sum_{k=1}^{j-2} \binom{n}{k} Q_{n-k,j}(t)
  +
  t \binom{n}{j-1} Q_{n-j+1,j}(t),
\end{equation}
where the middle sum is empty if $j=2$.
\end{proposition}
\begin{proof}
Let
\[
  G_j(z,t)
  \coloneqq
  \sum_{n \geq 0} Q_{n,j}(t)\,\frac{z^n}{n!}
  =
  t \sum_{n \geq 0} P_{n,j}(t)\,\frac{z^n}{n!}.
\]
By~\eqref{eq:big_block_egf},
\[
  G_j(z,t)
  =
  t\exp\!\left(
    \sum_{r=1}^{j-1} \frac{z^r}{r!}
    +
    t \sum_{r=j}^{\infty} \frac{z^r}{r!}
  \right).
\]
Differentiating with respect to~$z$ gives
\[
  \frac{\partial G_j}{\partial z}
  =
  \left(
    \sum_{r=0}^{j-2} \frac{z^r}{r!}
    +
    t \sum_{r=j-1}^{\infty} \frac{z^r}{r!}
  \right) G_j.
\]
Differentiating with respect to~$t$ gives
\[
  t \frac{\partial G_j}{\partial t}
  =
  \left(
    1
    +
    t \sum_{r=j}^{\infty} \frac{z^r}{r!}
  \right) G_j.
\]
Subtracting the second identity from the first yields
\[
  \frac{\partial G_j}{\partial z}
  =
  t \frac{\partial G_j}{\partial t}
  +
  \left(
    \sum_{r=1}^{j-2} \frac{z^r}{r!}
    +
    t\,\frac{z^{j-1}}{(j-1)!}
  \right) G_j.
\]
Extracting the coefficient of $z^n/n!$ gives~\eqref{eq:Q_recur}.
\end{proof}

\begin{conjecture}\label{conj:big_blocks_rr}
For every fixed $j \geq 2$, the polynomials $P_{n,j}(t)$ are real-rooted.
Moreover, they form a Sturm sequence:
\[
  P_{n,j}(t) \interl P_{n+1,j}(t)
  \qquad (n \geq 1).
\]
\end{conjecture}

\begin{remark}
If one prefers to count blocks of size \emph{strictly greater than}~$j$,
the same discussion applies after shifting the parameter from $j$ to $j+1$.
\end{remark}

\begin{remark}
Proposition~\ref{prop:touchard_rook_bijection} shows that the statistic
$\operatorname{bb}_j$ already has a natural staircase-rook interpretation:
it is the number of maximal rook chains with at least $j-1$ rooks in the
associated directed graph~$D(\rho)$. The remaining difficulty is not finding a
rook model, but proving the real-rootedness and interlacing conjectures for the
resulting generating polynomials.
\end{remark}

\subsection{Ordered set partitions}

Let $\Pi_n^{\operatorname{ord}}$ denote the set of ordered set partitions
of~$[n]$, and define
\[
  O_{n,j}(t)
  \coloneqq
  \sum_{\pi \in \Pi_n^{\operatorname{ord}}}
  t^{\operatorname{bb}_j(\pi)}.
\]
Since an ordered set partition is a finite sequence of nonempty blocks,
the symbolic method gives the exponential generating function
\begin{equation}\label{eq:ordered_big_block_egf}
  \sum_{n \geq 0} O_{n,j}(t)\,\frac{z^n}{n!}
  =
  \frac{1}{
    1
    -
    \left(
      \sum_{r=1}^{j-1} \frac{z^r}{r!}
      +
      t \sum_{r=j}^{\infty} \frac{z^r}{r!}
    \right)
  }.
\end{equation}

\begin{conjecture}\label{conj:ordered_big_blocks_rr}
For every fixed $j \geq 2$, the polynomials $O_{n,j}(t)$ are real-rooted.
Moreover, they form a Sturm sequence:
\[
  O_{n,j}(t) \interl O_{n+1,j}(t)
  \qquad (n \geq 0).
\]
\end{conjecture}

Computations currently verify Conjecture~\ref{conj:ordered_big_blocks_rr}
for $2 \leq j \leq 6$ and $n \leq 18$.
In this tested range, the ordered family looks even more rigid than the
unordered one: both real-rootedness and consecutive interlacing hold
throughout.

\subsection{Br\"and\'en--Leite and the generating functions}

The following result from \cite[Thm.~4.4]{BrandenLeite2024x}
gives real-rootedness and interlacing directly from the generating function.

\begin{theorem}[Br\"and\'en--Leite]\label{thm:brandenleite}
Let $f(x)$ be a polynomial with nonnegative coefficients
and all real zeros.
Consider the formal power series
\begin{equation}\label{eq:BLgf}
\frac{1}{1-t(f(x)-f(0))} = \sum_{n\geq 0} r_n(t)\, x^n \in \setR[t][[x]].
\end{equation}
Then $r_n(t)$ is real-rooted for each $n\in\setN$,
and $r_n(t) \interl r_{n+1}(t)$.
Moreover, if $f(0)\neq 0$, then all zeros of $r_n(t)$ lie in $[-1/f(0),\, 0]$.
\end{theorem}

\begin{theorem}[Br\"and\'en--Leite]\label{thm:brandenleite2}
Let $Q(x)$, where $Q(0)>0$, be a polynomial whose zeros are all real
and positive, and let $r$ be a positive integer.
Consider the power series
\begin{equation}\label{eq:BLgf2}
\frac{1}{Q(x)-tx^r} = \sum_{n\geq 0} q_n(t)\, x^n.
\end{equation}
Then $q_n(t)$ is a polynomial whose zeros are real and negative for each $n\in\setN$.
Moreover, $q_n(t) \interl q_{n+1}(t)$ for each $n\in\setN$.
\end{theorem}

\begin{remark}\label{rem:BL_not_direct}
Theorems~\ref{thm:brandenleite} and~\ref{thm:brandenleite2}
are tantalizingly close to the generating functions
\eqref{eq:big_block_egf} and~\eqref{eq:ordered_big_block_egf},
but do not apply directly.
For the unordered family, the exponential generating function
\eqref{eq:big_block_egf} is exponential rather than reciprocal.
For the ordered family, the denominator in~\eqref{eq:ordered_big_block_egf}
contains both a $t$-free part
\[
  \sum_{r=1}^{j-1} \frac{z^r}{r!}
\]
and a $t$-weighted tail
\[
  t \sum_{r=j}^{\infty} \frac{z^r}{r!},
\]
so it is neither of the form $1-t(f(z)-f(0))$ nor of the form
$Q(z)-t z^r$.
Thus Br\"and\'en--Leite does not currently give a direct proof for these
big-block polynomials, although it strongly suggests that some suitable
normalization or extension of the theorem might.
\end{remark}

\subsection{A Ferrers path-component extension}

The staircase construction extends naturally to arbitrary Ferrers boards.
Let $\mu = (\mu_1,\dotsc,\mu_\ell)$ be a partition and set
\[
  N(\mu) \coloneqq \max_{1 \leq i \leq \ell} (i + \mu_i).
\]
Embed the Ferrers board in the rectangle
$[\ell] \times [N(\mu)-1]$, reverse the columns, and interpret a cell
$(i,c)$ of~$\mu$ as the directed edge
\[
  i \longrightarrow N(\mu) - c + 1
\]
on the vertex set $[N(\mu)]$.
If $\rho$ is a standard rook placement on~$\mu$, let $D_\mu(\rho)$ be the
resulting directed graph.
Since $c \leq \mu_i$, every edge satisfies
\[
  N(\mu) - c + 1 \geq N(\mu) - \mu_i + 1 > i,
\]
so $D_\mu(\rho)$ is acyclic.
Because $\rho$ has at most one rook in each row and column, every vertex of
$D_\mu(\rho)$ has indegree and outdegree at most~$1$. Thus every connected
component of $D_\mu(\rho)$ is a directed path.

For $j \geq 2$, define
\[
  \operatorname{rch}^{\mu}_j(\rho)
  \coloneqq
  \#\{\text{path components of $D_\mu(\rho)$ with at least $j$ vertices}\},
\]
and let
\[
  C_{\mu,j}(t)
  \coloneqq
  \sum_{\rho} t^{\operatorname{rch}^{\mu}_j(\rho)},
\]
where the sum ranges over all standard rook placements~$\rho$ on~$\mu$.
For the staircase $\delta_{n-1}$ we have
\[
  N(\delta_{n-1}) = n,
\]
and Proposition~\ref{prop:touchard_rook_bijection} shows that
\[
  C_{\delta_{n-1},j}(t) = P_{n,j}(t).
\]

\begin{conjecture}\label{conj:ferrers_path_rr}
For every fixed $j \geq 2$ and every Ferrers board~$\mu$,
the polynomial $C_{\mu,j}(t)$ is real-rooted.
\end{conjecture}

Computations currently support Conjecture~\ref{conj:ferrers_path_rr}
for all $\mu$ with $|\mu| \leq 20$ when $j=2$, and for all $\mu$
with $|\mu| \leq 18$ when $j=3,4$.
Unlike the staircase case, however, the family
$\{C_{\mu,j}(t)\}$ does not appear to form a naive Sturm sequence under
row deletion.

\begin{remark}\label{rem:square_path}
For the square board $\square_n \coloneqq (n^n)$ one has
\[
  N(\square_n) = 2n.
\]
Hence every edge of $D_{\square_n}(\rho)$ goes from $[n]$ to
$\{n+1,\dotsc,2n\}$, so every path component has at most two vertices.
Therefore, for $j=2$ the statistic $\operatorname{rch}^{\square_n}_2(\rho)$
is simply the number of rooks in~$\rho$, and
\[
  C_{\square_n,2}(t)
  =
  \sum_{k=0}^{n} \binom{n}{k}^2 k!\, t^k
  =
  n! \, t^n L_n(-t^{-1}),
\]
where $L_n$ is the $n$-th Laguerre polynomial.
For $j \geq 3$, the polynomial $C_{\square_n,j}(t)$ is constant.
\end{remark}

\section{Further observations}

\subsection{Poset and LGV reformulation}

Define the \defin{strict Ferrers-cell poset}
\[
  P_\mu
  \coloneqq
  \{(i,j) : 1 \leq i \leq \ell,\ 1 \leq j \leq \mu_i\},
\]
ordered by
\[
  (i,j) < (i',j')
  \quad\Longleftrightarrow\quad
  i < i' \text{ and } j > j'.
\]

\begin{proposition}\label{prop:poset}
A non-nesting rook placement on~$\mu$ is exactly a chain in~$P_\mu$.
In particular,
\[
  R_\mu(t)
  = \sum_{k \geq 0} c_k(P_\mu)\, t^k,
\]
where $c_k(P_\mu)$ is the number of $k$-element chains in~$P_\mu$.
\end{proposition}
\begin{proof}
If $(i,j)$ and $(i',j')$ lie in a chain, then necessarily $i \neq i'$ and
$j \neq j'$, so a chain contains at most one cell from each row and column.
Moreover, after ordering the cells by increasing row index,
the column indices are strictly decreasing.
This is exactly the non-nesting condition.
Conversely, a non-nesting rook placement has strictly increasing row indices
and strictly decreasing column indices, so its cells form a chain in~$P_\mu$.
\end{proof}

For a fixed set of rows
$I = \{i_1 < \dotsb < i_k\} \subseteq [\ell]$,
let $d_I(\mu)$ be the number of non-nesting rook placements using exactly
those rows.

\begin{proposition}[LGV determinant for fixed rows]\label{prop:lgv}
For $I = \{i_1 < \dotsb < i_k\}$,
\begin{equation}\label{eq:lgv}
  d_I(\mu)
  =
  \det \!\left[
    \binom{\mu_{i_{k+1-b}} - b + 1}{a - b + 1}
  \right]_{a,b=1}^k.
\end{equation}
Consequently,
\begin{equation}\label{eq:rk_lgv}
  r_k(\mu)
  =
  \sum_{1 \leq i_1 < \dotsb < i_k \leq \ell}
  \det \!\left[
    \binom{\mu_{i_{k+1-b}} - b + 1}{a - b + 1}
  \right]_{a,b=1}^k.
\end{equation}
\end{proposition}
\begin{proof}
A placement on the rows $i_1 < \dotsb < i_k$ is determined by columns
$c_1 > \dotsb > c_k$ with $1 \leq c_a \leq \mu_{i_a}$.
Reversing the order by setting $x_b = c_{k+1-b}$ gives
\[
  1 \leq x_1 < x_2 < \dotsb < x_k
  \qquad\text{and}\qquad
  x_b \leq \mu_{i_{k+1-b}}.
\]
The number of such flagged strict subsets is given by the standard
Lindstr\"om--Gessel--Viennot determinant
\cite{Lindstrom1973,GesselViennot1985}, namely~\eqref{eq:lgv}.
Summing over all $k$-subsets of rows gives~\eqref{eq:rk_lgv}.
\end{proof}

\begin{remark}\label{rem:lgv}
Proposition~\ref{prop:lgv} packages each fixed-row contribution as a
flagged binomial determinant, hence as a local total-nonnegativity object.
The missing step is global: we do not yet see a single planar network
or a single totally nonnegative matrix whose relevant minors sum to
$r_k(\mu)$ for all~$k$.
\end{remark}

\subsection{Monotonicity under containment}

Conjecture~\ref{conj:main_strong} would immediately imply:
if $\nu$ is obtained from $\mu$ by stripping columns
(i.e., $\nu = \mu_{\langle c \rangle}$ for some $c \geq 1$),
then $R_\nu \interl R_\mu$.
A computational check for all pairs $\nu \subseteq \mu$
with $|\mu| \leq 14$ suggests the stronger conjecture that
$R_\nu \interl R_\mu$ whenever $\nu_i \leq \mu_i$ for all $i$
(containment of Ferrers boards).


%%%%%%%%%%%%%%%%%%%%%%%%%%%%%%%%%%%%%%%%%%%%%%%%%%%%%%
\section{Revised proof strategy}\label{sec:revised}
%%%%%%%%%%%%%%%%%%%%%%%%%%%%%%%%%%%%%%%%%%%%%%%%%%%%%%

The Br\"and\'en-style argument above has a gap:
the sequence $(f_1, \dotsc, f_{m+1})$
is only \emph{adjacent}-interlacing (from the column-strip of~$\mu'$),
not pairwise-interlacing, so Br\"and\'en's
Corollary~8.7 does not apply.
The rest of this section explains a revised strategy,
using the shift lemma and Wagner's cone lemma in place of
Br\"and\'en's matrix criterion.

\subsection{Tools}

We collect the lemmas we need; all are from
Wagner~\cite{Wagner1992},
Borcea--Br\"and\'en~\cite{BorceaBranden2010},
and Fisk~\cite{Fisk2006x}.
Throughout, all polynomials have positive leading coefficients
and only real non-positive roots.

\begin{lemma}[Wagner's cone lemma]\label{lem:wagner}
Let $f, g, h$ be real-rooted with non-positive roots.
\begin{enumerate}
\item If $f \interl h$ and $g \interl h$, then $f + g \interl h$.
\item If $h \interl f$ and $h \interl g$, then $h \interl f + g$.
\item $g \interl f$ if and only if $f \interl tg$.
\end{enumerate}
\end{lemma}

\begin{corollary}[Standard additive step]\label{cor:add_tf}
If $f \interl g$, then
\[
  g \interl g + t f.
\]
\end{corollary}
\begin{proof}
By Lemma~\ref{lem:wagner}(3), the hypothesis $f \interl g$ is equivalent to
$g \interl t f$.
Applying Lemma~\ref{lem:wagner}(2) to the trivial relation
$g \interl g$ and to $g \interl t f$ gives
$g \interl g + t f$.
\end{proof}

\begin{lemma}[Shift lemma]\label{lem:shift}
Let $f$ and $h$ be real-rooted with non-positive roots
and positive leading coefficients.
Suppose $h \interl f$ and $f(0) \geq h(0)$.
Then $g \coloneqq f + (t-1)h$ is real-rooted with
non-positive roots, and $f \interl g$.
\end{lemma}

\begin{proof}
At each root $r_i$ of $f$, $g(r_i) = (r_i - 1)\,h(r_i)$.
Since $r_i \leq 0$, the factor $r_i - 1$ is negative.
The interlacing $h \interl f$ forces $\operatorname{sign} h(r_i) = (-1)^{i-1}$,
so $g(r_i)$ alternates in sign across consecutive roots of~$f$.
The intermediate value theorem yields the required roots.
The boundary conditions (at~$0$ and at $-\infty$) follow from
$g(0) = f(0) - h(0) \geq 0$ and a leading-term comparison;
see~\cite{AlexanderssonJaldevik2025x} for the complete argument.
\end{proof}


\subsection{Notation and key identity}

As before, set $m = \mu_\ell$, $\mu' = (\mu_1, \dotsc, \mu_{\ell-1})$,
and $f_c = R_{\mu'_{\langle c \rangle}}$.
By the outer induction hypothesis,
\begin{equation}\label{eq:outer_ih}
  f_{c+1} \interl f_c \quad\text{for all } 0 \leq c < \mu_1.
\end{equation}

For $0 \leq c \leq s \leq m$, define the \defin{partial sums}
\[
  G(c, s)
  \coloneqq f_c + t \sum_{j=c+1}^{s} f_j,
  \qquad G(c, c) = f_c.
\]
The recursion gives $R_{\mu_{\langle c \rangle}} = G(c, m)$ for $0 \leq c < m$.

\begin{observation}[Constant difference]\label{obs:const_diff}
For all $s$,
\[
  G(c, s) - G(c{+}1, s) = f_c + (t-1)\,f_{c+1}
  \eqqcolon D_c.
\]
In particular, $D_c$ is independent of~$s$.
\end{observation}

Since $f_c(0) = f_{c+1}(0) = 1$, we have $D_c(0) = 0$, so $D_c = t \cdot (D_c / t)$.
We define
\[
  \defin{\psi_c} \coloneqq \frac{f_c - f_{c+1}}{t}
  \qquad\text{(well-defined since $f_c(0) = f_{c+1}(0)$)}.
\]
Then $D_c / t = f_{c+1} + \psi_c$.

\subsection{Difference quotients and a grid package}\label{sec:grid_package}

For any Ferrers board~$\nu$ and $0 \leq c \leq \nu_1$, define
\[
  F_c^\nu \coloneqq R_{\nu_{\langle c \rangle}},
\]
For $0 \leq c < \nu_1$, set
\[
  \Delta_c^\nu \coloneqq \frac{F_c^\nu - F_{c+1}^\nu}{t}.
\]
Since $F_c^\nu(0) = F_{c+1}^\nu(0) = 1$, the quotient $\Delta_c^\nu$
is a polynomial with integer coefficients.

Now let $\mu = (\mu', m)$ with $m = \mu_\ell$, and keep the notation
$f_c = F_c^{\mu'}$.
For $0 \leq c \leq s \leq m$, recall
\[
  G(c,s) = f_c + t \sum_{j=c+1}^{s} f_j,
\]
so that $R_{\mu_{\langle c \rangle}} = G(c,m)$.

\begin{lemma}[Difference-quotient recurrences]\label{lem:delta_grid}
For every $0 \leq c < m$ and $c+1 \leq s \leq m$,
\begin{align}
  \Delta_c^\mu
    &= \Delta_c^{\mu'} + F_{c+1}^{\mu'}
     = \frac{f_c - f_{c+1}}{t} + f_{c+1}, \label{eq:delta_recur} \\
  G(c,s)
    &= G(c+1,s) + t\,\Delta_c^\mu, \label{eq:grid_horizontal} \\
  G(c,s+1)
    &= G(c,s) + t\,F_{s+1}^{\mu'}. \label{eq:grid_vertical}
\end{align}
\end{lemma}
\begin{proof}
The identity~\eqref{eq:delta_recur} is just
\[
  \Delta_c^\mu
  = \frac{R_{\mu_{\langle c \rangle}} - R_{\mu_{\langle c+1 \rangle}}}{t}
  = \frac{G(c,m) - G(c+1,m)}{t}
  = \frac{f_c - f_{c+1}}{t} + f_{c+1},
\]
using Observation~\ref{obs:const_diff}.
Equation~\eqref{eq:grid_horizontal} is the identity
$G(c,s)-G(c+1,s) = D_c = t\,\Delta_c^\mu$.
Equation~\eqref{eq:grid_vertical} is immediate from the definition of~$G(c,s)$.
\end{proof}

The horizontal recurrence~\eqref{eq:grid_horizontal} isolates the key
auxiliary polynomial $\Delta_c^\mu$.
Indeed, if one could prove
\begin{equation}\label{eq:delta_suffices}
  \Delta_c^\mu \interl G(c+1,s)
  \qquad (0 \leq c < s \leq m),
\end{equation}
then adjacent strip interlacing would follow immediately:
by Wagner's lemma,
$\Delta_c^\mu \interl G(c+1,s)$ is equivalent to
$G(c+1,s) \interl t\,\Delta_c^\mu$,
and the cone then gives
\[
  G(c+1,s)
  \interl G(c+1,s) + t\,\Delta_c^\mu
  = G(c,s).
\]
In particular, taking $s=m$ yields
$R_{\mu_{\langle c+1 \rangle}} \interl R_{\mu_{\langle c \rangle}}$.

\begin{remark}[Computational grid package]\label{rem:grid_package}
The binary \texttt{nn\_rook\_grid.rs}
in \texttt{experiments/src/bin/}
checks the recurrence package above for all partitions with $|\mu| \leq 20$.
Every test passes:
\begin{itemize}
\item $\Delta_c^{\mu'} = (f_c - f_{c+1})/t$ is real-rooted and coefficientwise nonnegative:
  $17605/17605$.
\item $\Delta_c^\mu$ is real-rooted and coefficientwise nonnegative:
  $3738/3738$.
\item $\Delta_c^\mu \interl f_{c+1}$:
  $3738/3738$.
\item $\Delta_c^\mu \interl G(c+1,s)$ and $\Delta_c^\mu \interl G(c,s)$:
  $5657/5657$ in each case.
\item Hence $G(c+1,s) \interl G(c,s)$:
  $5657/5657$.
\item In addition, the diagonal relations
  $G(c+1,s) \interl G(c,s+1)$ and
  $G(c+1,s+1) \interl G(c,s)$
  both hold:
  $1919/1919$ in each case.
\end{itemize}
Computationally, this $\Delta$-package is substantially cleaner than the
earlier $\psi_c$-based formulation below, and seems the right recurrence
system to target in a proof.
\end{remark}

\subsection{Interlacing squares}\label{sec:squares}

Fix $0 \leq c < s < m$ and write
\[
  A \coloneqq G(c+1,s),
  \qquad
  B \coloneqq G(c,s) = A + t\,\Delta_c^\mu,
  \qquad
  R \coloneqq F_{s+1}^{\mu'}.
\]
Then the corresponding $2 \times 2$ square in the grid is
\[
  \begin{matrix}
    A & \interl? & B \\
    \interl? & & \interl? \\
    A + tR & \interl? & B + tR
  \end{matrix}
  \qquad
  \text{where }
  A + tR = G(c+1,s+1),\ B + tR = G(c,s+1).
\]

By Corollary~\ref{cor:add_tf}, the vertical relations
\[
  A \interl A+tR
  \qquad\text{and}\qquad
  B \interl B+tR
\]
would follow from $R \interl A$ and $R \interl B$, respectively.
Computationally, however, these hypotheses fail frequently.
The new binary \texttt{nn\_rook\_square.rs}
checks all such squares for $|\mu| \leq 14$ and finds:
\begin{itemize}
\item the row relations
  $A \interl B$ and $(A+tR) \interl (B+tR)$
  both hold in all $333/333$ squares;
\item the left vertical relation $A \interl A+tR$ holds in $290/333$ squares;
\item the right vertical relation $B \interl B+tR$ holds in $195/333$ squares;
\item nevertheless, both the diagonal relation
  $A \interl B+tR$ and the anti-diagonal relation
  $A+tR \interl B$
  hold in all $333/333$ squares.
\end{itemize}

The diagonal relation has a clean experimental explanation.
Indeed,
\[
  B+tR
  = A + t\,(\Delta_c^\mu + R),
\]
and the same computation verifies
\begin{equation}\label{eq:diag_candidate}
  \Delta_c^\mu + F_{s+1}^{\mu'} \interl G(c+1,s)
  \qquad\text{in all } 333/333 \text{ tested squares.}
\end{equation}
Thus Corollary~\ref{cor:add_tf} would imply the diagonal edge
\[
  G(c+1,s) \interl G(c,s+1).
\]

The anti-diagonal admits a parallel reformulation.
Since
\[
  B
  = (A+tR) + t\,(\Delta_c^\mu - R),
\]
the anti-diagonal edge
\[
  G(c+1,s+1) \interl G(c,s).
\]
would follow from
\begin{equation}\label{eq:antidiag_candidate}
  \Delta_c^\mu - F_{s+1}^{\mu'}
  \interl G(c+1,s+1).
\end{equation}
The same square computation verifies two stronger facts in all
$333/333$ tested squares:
\begin{itemize}
\item $\Delta_c^\mu - F_{s+1}^{\mu'}$ has nonnegative coefficients;
\item $\Delta_c^\mu - F_{s+1}^{\mu'} \interl G(c+1,s+1)$ and
  $\Delta_c^\mu - F_{s+1}^{\mu'} \interl G(c,s)$.
\end{itemize}

Thus the square seems to be governed by two coupled difference quotients:
\[
  \Delta_c^\mu + F_{s+1}^{\mu'}
  = \frac{G(c,s+1)-G(c+1,s)}{t},
  \qquad
  \Delta_c^\mu - F_{s+1}^{\mu'}
  = \frac{G(c,s)-G(c+1,s+1)}{t}.
\]
At present we do not know a conceptual proof of
\eqref{eq:diag_candidate} or \eqref{eq:antidiag_candidate},
but these diagonal and anti-diagonal phenomena seem more robust than the
vertical edges, and may be the right local structure to target.

\subsection{Verified properties for the \texorpdfstring{$\psi_c$}{psi_c} package}\label{sec:properties}

The following six properties have been verified computationally
for all partitions~$\mu$ with $|\mu| \leq 14$
(see \texttt{experiments/src/bin/nn\_rook\_proof\_final.rs}).
Properties~(P1)--(P3) are used in the proof;
(P4)--(P6) record additional structure.

\begin{enumerate}[label=(P\arabic*)]
\item\label{P:ih} \textbf{Outer IH.}
  $f_{c+1} \interl f_c$ for all $0 \leq c < \mu_1$.
  \hfill{(induction hypothesis)}

\item\label{P:psi_nonneg} \textbf{$\psi_c$ has nonnegative coefficients}
  for all $0 \leq c < m$.
  \hfill{(693/693)}

\item\label{P:D_interl_f} \textbf{$D_c/t \interl f_{c+1}$}
  for all $0 \leq c < m$, i.e., $D_c/t$ has roots
  strictly to the left of those of~$f_{c+1}$.
  \hfill{(693/693)}

\item\label{P:D_interl_G} \textbf{$D_c/t \interl G(c{+}1, s)$}
  for all $0 \leq c < m$ and $c{+}1 \leq s \leq m$.
  \hfill{(1026/1026)}

\item\label{P:main} \textbf{$G(c{+}1, s) \interl G(c, s)$}
  for all $0 \leq c < m$ and $c{+}1 \leq s \leq m$.
  \hfill{(1026/1026)}

\item\label{P:psi_rr} \textbf{$\psi_c$ is real-rooted}
  for all $0 \leq c < m$.
  \hfill{(693/693)}
\end{enumerate}

\subsection{How Property (P4) Would Imply Adjacent Interlacing}

The adjacent column-strip interlacing
$R_{\mu_{\langle c+1 \rangle}} \interl R_{\mu_{\langle c \rangle}}$
is property~\ref{P:main} at $s = m$.
The next lemma shows that the missing ingredient is precisely
property~\ref{P:D_interl_G}.

\begin{lemma}[From \ref{P:D_interl_G} to \ref{P:main}]\label{lem:P4_implies_P5}
If~\ref{P:D_interl_G} holds, then~\ref{P:main} holds.
\end{lemma}
\begin{proof}
Fix $c$ and $s$.
Since $D_c/t \interl G(c{+}1, s)$,
Wagner's lemma (part~3) gives
$G(c{+}1, s) \interl t \cdot (D_c/t) = D_c$.
The left cone (part~2 of Wagner) applied to
$G(c{+}1, s) \interl G(c{+}1, s)$ (trivially)
and $G(c{+}1, s) \interl D_c$ yields
\[
  G(c{+}1, s) \interl G(c{+}1, s) + D_c = G(c, s). \qedhere
\]
\end{proof}

\begin{lemma}[Shift lemma gives $D_c/t \interl f_c$]\label{lem:shift_gives_Dc}
Property~\ref{P:ih} implies $D_c/t \interl f_c$.
\end{lemma}
\begin{proof}
By~\ref{P:ih}, $f_{c+1} \interl f_c$ and $f_c(0) = 1 \geq 1 = f_{c+1}(0)$.
The shift lemma gives $f_c \interl f_c + (t{-}1)f_{c+1} = D_c = t \cdot (D_c/t)$.
By Wagner (part~3), this is equivalent to $D_c/t \interl f_c$.
\end{proof}

It remains to prove~\ref{P:D_interl_G}.
The base case ($s = c{+}1$, i.e., $G(c{+}1, c{+}1) = f_{c+1}$)
is exactly~\ref{P:D_interl_f}.
The inductive step
($D_c/t \interl G(c{+}1, s)
 \Rightarrow D_c/t \interl G(c{+}1, s{+}1) = G(c{+}1, s) + t f_{s+1}$)
would follow from the left cone if
$D_c/t \interl t f_{s+1}$,
equivalently $f_{s+1} \interl D_c/t$ by Wagner.
However, this last condition \textbf{fails} in general
(1508/2598 in the computational check).

\begin{openproblem}\label{open:key}
Find a proof of~\ref{P:D_interl_G}
(or equivalently~\ref{P:main})
that does not require $f_{s+1} \interl D_c/t$.
In particular:
\begin{itemize}
\item Does the $q$-deformation $R_\mu(q, t) = \sum q^{\mathrm{nest}} t^k$
  (which is real-rooted in~$t$ for all $q \in [0,1]$)
  provide additional structure?
\item Can~\ref{P:D_interl_G} be proved directly
  from~\ref{P:D_interl_f} (the base case) and
  the identity $G(c{+}1, s{+}1) = G(c{+}1, s) + t f_{s+1}$
  by a continuous deformation argument?
  (Computationally, interlacing is maintained along
  $\alpha \mapsto G(c{+}1, s) + \alpha\, t f_{s+1}$
  for $\alpha \in [0, 1]$; verified at 101 points
  for $|\mu| \leq 14$.)
\end{itemize}
\end{openproblem}

\subsection{What is proved unconditionally}

Summarizing:
\begin{itemize}
\item The difference-quotient identities
  \eqref{eq:delta_recur}--\eqref{eq:grid_vertical}
  reduce adjacent strip interlacing to the single missing statement
  $\Delta_c^\mu \interl G(c{+}1,s)$, and this entire grid package
  is verified computationally for $|\mu| \leq 20$
  (Remark~\ref{rem:grid_package}).
\item The shift lemma gives
  $D_c/t \interl f_c$ (Lemma~\ref{lem:shift_gives_Dc}).
\item \ref{P:D_interl_G}${}\Rightarrow{}$\ref{P:main}
  via Wagner's cone (Lemma~\ref{lem:P4_implies_P5}).
\item The base case~\ref{P:D_interl_f} of~\ref{P:D_interl_G}
  is strictly stronger than $D_c/t \interl f_c$
  and is verified but not yet proved from~\ref{P:ih} alone.
\item Properties~\ref{P:psi_nonneg} and~\ref{P:psi_rr}
  ($\psi_c$ nonneg and real-rooted) hold computationally
  and may provide the missing ingredient.
\end{itemize}


\begin{thebibliography}{99}

\bibitem{AlexanderssonJaldevik2025x}
P.~Alexandersson, O.~Jaldevik,
\emph{Rook-Eulerian polynomials and permutation ideals},
preprint, 2025.

\bibitem{BorceaBranden2010}
J.~Borcea, P.~Br\"and\'en,
\emph{Multivariate P\'olya--Schur classification problems
in the Weyl algebra},
Proc.\ London Math.\ Soc., 101 (2010), 73--104.

\bibitem{Branden2015}
P.~Br\"and\'en,
\emph{Unimodality, log-concavity, real-rootedness and beyond},
Handbook of Enumerative Combinatorics, 2015.

\bibitem{BrandenLeite2024x}
P.~Br\"and\'en, L.~S.~M.~Leite,
\emph{Totally nonnegative matrices, chain enumeration and zeros of polynomials},
arXiv:2412.06595, 2024.

\bibitem{Fisk2006x}
S.~Fisk,
\emph{Polynomials, roots, and interlacing},
preprint, 2006.

\bibitem{GesselViennot1985}
I.~Gessel, G.~Viennot,
\emph{Binomial determinants, paths, and hook length formulae},
Adv.\ Math., 58 (1985), 300--321.

\bibitem{Lindstrom1973}
B.~Lindstr\"om,
\emph{On the vector representations of induced matroids},
Bull.\ London Math.\ Soc., 5 (1973), 85--90.

\bibitem{Wagner1992}
D.~Wagner,
\emph{Total positivity of Hadamard products},
J.~Math.~Anal.~Appl., 163 (1992), 459--483.
\end{thebibliography}

\end{document}
